% !TeX encoding = UTF-8
We sincerely thank the referee for the careful reading of our manuscript and for the constructive report. The referee's comments helped us improve the motivation, structure, clarity, and practical usefulness of the paper. Below, we respond to each part of the report in the same order as the referee's comments: Summary of the work, Recommendation, Major points, and Minor points.
\section*{Comment 1: Summary of the work}
\begin{commentbox}
\textbf{Summary of the work.} This manuscript studies whether structured note-taking can help students understand difficult reading materials. The authors compare ordinary linear notes with a more organized note format that asks students to record the main idea, key terms, examples, and questions on the same page. The topic is timely and potentially useful for students and teachers. The manuscript is generally clear, and the classroom example is easy to follow. However, the paper should explain more clearly what is new about the proposed note-taking format and how it differs from existing methods such as mind maps, Cornell notes, or ordinary outlines.
\end{commentbox}
\begin{replybox}
We thank the referee for the positive and balanced summary of the work. We are pleased that the referee found the topic timely and the classroom example easy to follow.We also agree that the original manuscript did not state the novelty of the proposed note-taking format clearly enough. In the revised manuscript, we have added a new paragraph at the end of the Introduction to explain the difference between our structured note sheet and commonly used note-taking methods. In particular, we now clarify that our method is not intended to replace mind maps, Cornell notes, or outlines. Instead, it combines three simple elements in one page: a short summary, a list of key terms, and a question area for unclear points.To make the distinction more concrete, we added a comparison table in Section 2. The table explains that mind maps are especially useful for showing relations, Cornell notes are useful for review, outlines are useful for hierarchy, while our format is designed for first-time reading of difficult texts. This revision should help readers understand the specific contribution of the paper.
\end{replybox}
\section*{Comment 2: Recommendation}
\begin{commentbox}
\textbf{Recommendation.} I recommend that the manuscript be considered for publication after major revision. The topic is valuable and the paper has practical potential, but the authors should revise the manuscript substantially before acceptance. In particular, the authors should clarify the novelty, describe the classroom procedure in more detail, improve the explanation of the data, and discuss the limitations more carefully.
\end{commentbox}
\begin{replybox}
We thank the referee for the recommendation and for identifying the main areas requiring revision. We have treated the report as a guide for a substantial revision rather than as a list of minor corrections.In response to the recommendation, we have made four main changes. First, we clarified the novelty of the proposed note-taking format in the Introduction and Section 2. Second, we expanded the description of the classroom procedure, including the reading task, time arrangement, student instructions, and scoring method. Third, we revised the Results section to explain the data in plain language before presenting numerical comparisons. Fourth, we added a more explicit Limitations subsection, discussing sample size, course context, and possible differences among students with different reading habits.We believe that these revisions address the referee's main concerns and make the manuscript more suitable for publication.
\end{replybox}
\section*{Major points}
\section*{Comment 3: Novelty of the proposed method}
\begin{commentbox}
\textbf{Major point 1.} The manuscript should explain more clearly what is new about the proposed note-taking method. At present, the method appears similar to existing classroom tools, especially mind maps, Cornell notes, and outlines.
\end{commentbox}
\begin{replybox}
We agree with the referee. The original version introduced the note-taking sheet too quickly and did not sufficiently distinguish it from existing methods.We have revised Section 2 to explain the method more explicitly. The revised text now states that the proposed sheet is designed for students who are reading a difficult text for the first time. Its purpose is to prevent students from copying sentences mechanically. Therefore, the sheet requires students to complete four simple tasks while reading: identify the main idea, record key terms, write one example, and raise at least one question.We also added a table comparing the proposed method with mind maps, Cornell notes, and outlines. This table makes clear that our contribution lies in providing a low-threshold, one-page structure for first-time reading rather than proposing an entirely new theory of note-taking.
\end{replybox}
\section*{Comment 4: Description of the classroom procedure}
\begin{commentbox}
\textbf{Major point 2.} The classroom procedure needs more detail. The reader should know what the students were asked to do, how much time they had, and how the teacher introduced the note-taking sheet.
\end{commentbox}
\begin{replybox}
Thank you for this helpful suggestion. We have expanded the description of the classroom procedure in Section 3.The revised manuscript now explains the procedure step by step. Students first received a short article of about 900 words. They were then given 20 minutes to read the article and take notes. The control group used their usual note-taking style, while the experimental group used the structured note sheet. Before the task began, the teacher gave a five-minute demonstration showing how to fill in the four parts of the sheet.We also added the exact instruction given to students: ``Please do not copy the whole paragraph. Try to write the main idea in your own words, mark important terms, give one example, and write down one question.'' We believe this addition makes the study easier to understand and easier to reproduce.
\end{replybox}
\section*{Comment 5: Explanation of the data}
\begin{commentbox}
\textbf{Major point 3.} The Results section presents several numbers, but the meaning of these numbers is not sufficiently explained. Please make the data more readable for non-specialist readers.
\end{commentbox}
\begin{replybox}
We appreciate this comment. The original Results section was too compressed and assumed too much statistical familiarity from the reader.In the revised manuscript, we now introduce each result in plain language before giving the numerical values. For example, before reporting the score difference between the two groups, we first explain that the score measures how accurately students could explain the main idea of the reading passage after the task.We also revised Figure 2 and Table 1. The figure now uses clearer labels, and the table includes a short note explaining what each score means. In addition, we moved a technical sentence about the statistical test to a footnote so that the main text remains readable.
\end{replybox}
\section*{Comment 6: Practical significance}
\begin{commentbox}
\textbf{Major point 4.} The manuscript reports that the structured note-taking group performed better, but it should explain whether the difference is educationally meaningful. A small numerical difference may not always matter in practice.
\end{commentbox}
\begin{replybox}
We agree that statistical improvement alone is not enough. The paper should also explain whether the result matters in an actual classroom.We have added a paragraph in Section 4 to discuss practical significance. The revised manuscript now explains that the improvement is meaningful because students in the structured note-taking group not only received higher scores but also produced more complete explanations of the reading passage. In other words, the difference was visible not only in the number but also in the quality of the students' written responses.We also added two anonymized examples of student answers. These examples show that students using the structured sheet were more likely to connect the main idea with supporting details. This makes the practical value of the method clearer.
\end{replybox}
\section*{Comment 7: Limitations and generalizability}
\begin{commentbox}
\textbf{Major point 5.} The limitations are mentioned only briefly. The authors should discuss whether the findings can be generalized to other subjects, younger students, or longer reading tasks.
\end{commentbox}
\begin{replybox}
Thank you for this important comment. We have expanded the Limitations subsection substantially.The revised manuscript now discusses three limitations. First, the study involved undergraduate students from one university, so the findings should be tested in other educational contexts. Second, the reading task was relatively short; longer readings may require different note-taking strategies. Third, the study focused on expository texts, and the results may not directly apply to literary texts, laboratory reports, or mathematical proofs.We also added a short paragraph on future work. We now suggest testing the method with younger students, comparing different subject areas, and studying whether repeated use of the note sheet improves long-term reading habits.
\end{replybox}
\section*{Minor points}
\section*{Comment 8: Title}
\begin{commentbox}
\textbf{Minor point 1.} The title is understandable but slightly broad. Please consider making it more specific.
\end{commentbox}
\begin{replybox}
Thank you. We have revised the title from ``A Simple Way to Improve Learning'' to ``How Structured Note-Taking Helps Students Understand Difficult Texts.'' The new title states the method, the target activity, and the expected outcome more clearly.
\end{replybox}
\section*{Comment 9: Abstract}
\begin{commentbox}
\textbf{Minor point 2.} The abstract should mention the sample size and the main finding more directly.
\end{commentbox}
\begin{replybox}
We agree. The revised abstract now states that 120 undergraduate students participated in the study. It also states the main finding more directly: students using the structured note sheet showed better main-idea recall and produced more complete short explanations.
\end{replybox}
\section*{Comment 10: Definition of ``structured note-taking''}
\begin{commentbox}
\textbf{Minor point 3.} The term ``structured note-taking'' should be defined the first time it appears.
\end{commentbox}
\begin{replybox}
Thank you for pointing this out. We now define structured note-taking in the first paragraph of the Introduction as ``a note-taking approach that guides students to organize information into fixed parts, such as main ideas, key terms, examples, and questions.''
\end{replybox}
\section*{Comment 11: Figure caption}
\begin{commentbox}
\textbf{Minor point 4.} The caption of Figure 1 is too short. Please explain what the four boxes in the note sheet mean.
\end{commentbox}
\begin{replybox}
We have revised the caption of Figure 1. The new caption explains that the four boxes correspond to the main idea, key terms, supporting example, and reader's question. It also states that the sheet is intended to guide students during first-time reading.
\end{replybox}
\section*{Comment 12: Wording}
\begin{commentbox}
\textbf{Minor point 5.} The phrase ``students become smarter readers'' sounds informal. Please use more precise wording.
\end{commentbox}
\begin{replybox}
We agree. We have replaced this phrase with ``students develop more organized reading strategies.'' This wording is more precise and more appropriate for an academic manuscript.
\end{replybox}
\section*{Comment 13: Formatting and references}
\begin{commentbox}
\textbf{Minor point 6.} Please check the consistency of capitalization in section headings and the reference list.
\end{commentbox}
\begin{replybox}
Thank you. We have checked all section headings and references. Capitalization has been made consistent throughout the manuscript, and the reference list has been revised to follow the journal style.
\end{replybox}
\vspace{1em}
We once again thank the referee for the careful and constructive review. We believe that the revised manuscript is clearer, more practical, and more accessible to a broad readership.
